\chapter{Introduction}
\label{ch:intro}

\ccprod{} is a spreadsheet program for the Commander X16. It holds up to
sixteen worksheets in one workbook, sixteen thousand cells on each of them,
and recalculates a formula written in terms of any of the others. It reads and
writes its own file format, imports and exports comma-separated text, and
imports and exports Excel \fname{.XLSX} workbooks. It does all of that in a
thirty-kilobyte program on a 65C02 running at eight megahertz.

That last sentence is the reason for most of this manual. \ccprod{} is not a
scaled-down version of a large program; it is a small program, designed as one,
and the shape of it shows. Knowing what it does \emph{not} do is as useful as
knowing what it does, so this manual says both.

\section{What Commander Calc is}

\begin{ccsteps}
\item \textbf{A worksheet program.} A grid of 256 columns by 65,535 rows,
addressed \cellref{A1} to \cellref{IV65535}. Cells hold numbers, text, dates,
\msg{TRUE} and \msg{FALSE}, formulas, and error values.
\item \textbf{A workbook program.} Up to sixteen worksheets in one file, with
their own names and their own tabs, and formulas that reach across them:
\fml{=SUM(Q1!B2:B40)} does what it looks like.
\item \textbf{A block editor.} Mark a rectangle of any size, then copy it,
paste it somewhere else, or clear it in one command.
\item \textbf{Formulas that move with their cells.} A pasted formula has its
references rewritten by the distance the block travelled, and so does one
carried along by a sort or by an inserted row. A \fml{\$} pins whatever should
not move.
\item \textbf{Charts.} Three commands draw the numbers in the cursor's column
as bars, as a line, or as a pie, labelled from the column beside them --- and
any of them can be saved to the card as a picture file.
\item \textbf{A calculator with fourteen functions.} \fn{SUM}, \fn{AVERAGE},
\fn{MIN}, \fn{MAX}, \fn{COUNT}, \fn{IF}, \fn{AND}, \fn{OR}, \fn{NOT},
\fn{ABS}, \fn{ROUND}, \fn{INT}, \fn{MOD} and \fn{LEN}. That is the whole list;
Chapter~\ref{ch:funcref} documents every one.
\item \textbf{A program that reads other people's files.} An Excel workbook
saved as \fname{.XLSX} can be opened, with its numbers, text, dates, number
formats, column widths, multiple worksheets and --- as far as the fourteen
functions reach --- its formulas.
\item \textbf{A program that writes files other people can read.} A workbook
can be written back out as \fname{.XLSX} or as \fname{.CSV}.
\end{ccsteps}

\section{What Commander Calc is not}

It is not Lotus~1-2-3 and it is not Excel, and it does not pretend to be. The
list below is the short version; Appendix~\ref{app:differences} is the long one.

\begin{ccsteps}
\item \textbf{There is no cut or move.} Copy and paste, and delete what you
left behind.
\item \textbf{There are no colour or font commands.} \ccprod{} sets column
widths and freezes headings, and it displays eight number formats, bold and
three alignments faithfully --- but the only way to get a \emph{format} into a
workbook is to import a file that already has one. Nothing on the menus sets a
number format, bold, or a colour.
\item \textbf{There are no text functions, lookup functions or date
functions.} No \fn{LEFT}, no \fn{VLOOKUP}, no \fn{SUMIF}, no \fn{TODAY}.
\item \textbf{Undo is one level deep.} One cell. Sorting has its own separate
one-level undo, and a pasted block cannot be taken back at all.
\item \textbf{A formula longer than 128 characters is never rewritten.} It is
left exactly as it was rather than half-adjusted, wherever it is moved to.
\item \textbf{There are no macros, no printing and no windows}, and a chart
covers one column at a time.
\end{ccsteps}

\begin{ccnote}
Most of these are absences rather than defects: work that has not been done
rather than work that was done wrong. Where something is actually broken ---
and there are one or two --- this manual says so at the point where you would
meet it.
\end{ccnote}

\section{What you need}
\index{requirements}

\begin{cctable}{You need}{Notes}{1.5in}{3.32in}
A Commander X16 & Or the \fname{x16emu} emulator, which is what most of this
manual was written against. \\
512K of banked RAM & The machine's standard fitting. The workbook lives
entirely in banked RAM, so more of it means a bigger workbook; \ccprod{} will
use up to the 2M maximum if the machine has it. \\
A storage device at unit 8 & An SD card under CMDR-DOS, or the emulator's host
filesystem. \\
An 80 by 60 text display & The default. \ccprod{} adapts to a smaller text mode
and refuses to run in anything shorter than ten lines. \\
A mouse & Optional. Everything can be done from the keyboard; the mouse is
supported for menus, dialogues, moving the cursor and scrolling. \\
\end{cctable}

\section{Where to go next}

\begin{cctable}{If you want to}{Read}{1.85in}{2.97in}
Get the program running & Chapter~\ref{ch:installing} \\
Understand the display & Chapter~\ref{ch:screen} \\
Type something into a cell & Chapters~\ref{ch:moving} and \ref{ch:entering} \\
Save your work & Chapter~\ref{ch:files} \\
Open a spreadsheet from another program & Chapter~\ref{ch:importexport} \\
Look up one command & Chapter~\ref{ch:cmdref} \\
Look up one key & Chapter~\ref{ch:keyref} \\
Look up one function & Chapter~\ref{ch:funcref} \\
Draw a picture of some numbers & Chapter~\ref{ch:charting} \\
Understand a message on the screen & Chapter~\ref{ch:messages} \\
Know exactly how big a workbook can be & Appendix~\ref{app:limits} \\
Know what is missing and why & Appendix~\ref{app:differences} \\
\end{cctable}
